\BourbakiExerciseGroup

\begin{enumerate}
\def\labelenumi{\arabic{enumi})}
\item
  当 M 是实数加法群 R 的子群时，将一元多项式中的欧几里得除法推广到 M 的群代数上。
\item
  设 \(f\) 为一个多项式 \(\neq 0\) ，它属于 \(\mathbf{A}[\mathbf{X}]\) ，次数为 \(n\) ，首项系数为 \(\alpha_0\) 。设 \(\mathbf{M}\) 为 \(\mathbf{A}[\mathbf{X}]\) （视为 A-模）的子模，由这样的多项式组成：其 \(m\) 次项（ \(m\) 为 \(\mathbf{N}\) 中任意元）的系数可被 \(\alpha_0^\mu\) 整除，其中 \(\mu = \max(m - n + 1, 0)\) ；证明对任意多项式 \(g\) （在 \(\mathbf{M}\) 中），存在两个多项式 \(u, v\) （在 \(\mathbf{A}[\mathbf{X}]\) 中），使得 \(\deg v < n\) 且 \(g = uf + v\) 。
\item
  \begin{enumerate}
  \def\labelenumii{\alph{enumii})}
  \tightlist
  \item
    在 A 上的一元多项式代数 A{[}X{]} 中，映射 \((u,v)\mapsto u(v)\) 是一个内部合成法则；证明该法则满足结合律，并且关于 A{[}X{]} 中的加法与乘法是左分配的。
  \end{enumerate}
\end{enumerate}

\begin{enumerate}
\def\labelenumi{\alph{enumi})}
\setcounter{enumi}{1}
\item
  若 A 是整环，证明关系 \(u(v)=0\) 蕴含 u=0 或 v 为常数；并且当 \(u\neq0\) 且 \(\deg v>0\) 时， \(u(v)\) 的次数等于 u 与 v 的次数之积。
\item
  从现在起假设 A 是一个域。设 u 和 v 是次数 \textgreater0 的两个多项式（在 A{[}X{]} 中）， f 是次数 \textgreater0 的多项式；证明：若 q 是 u 除以 v 的欧几里得除法的商， r 是其余式，则 \(q(f)\) 与 \(r(f)\) 是 \(u(f)\) 除以 \(v(f)\) 的欧几里得除法的商与余式。
\item
  对每个多项式 \(f\) （属于 \(\mathbf{A}[\mathbf{X}]\) ），用 \(\mathrm{I}(f)\) 表示所有形如 \(u(f)\) 的多项式的集合，其中 \(u\) 取遍 \(\mathbf{A}[\mathbf{X}]\) ；它是 \(\mathbf{A}[\mathbf{X}]\) 的一个子环。为使 \(\mathrm{I}(f) = \mathrm{I}(g)\) ，必要且充分的是 \(g = \lambda f + \mu\) ，其中 \(\lambda \neq 0\) 与 \(\mu\) 均属于 \(\mathbf{A}\) 。
\item
  设 \(f\) 与 \(g\) 为两个次数 \(> 0\) 的多项式（在 \(\mathbf{A}[\mathbf{X}]\) 中）；证明 \(I(f) \cap I(g)\) 或者等于 \(\mathbf{A}\) ，或者具有形式 \(I(h)\) ，其中 \(h\) 是次数 \(> 0\) 的多项式（排除第一种可能性，在 \(I(f) \cap I(g)\) 中考虑一个多项式 \(h\) ，其次数 \(> 0\) 且尽可能低）。
\end{enumerate}

\begin{enumerate}
\def\labelenumi{\arabic{enumi})}
\setcounter{enumi}{3}
\tightlist
\item
  \begin{enumerate}
  \def\labelenumii{\alph{enumii})}
  \tightlist
  \item
    假设对每个非零的 \(n \in \mathbf{Z}\) ，关系 \(n\xi = 0\) 蕴含 \(\xi = 0\) （在 \(\mathbf{A}\) 中）。证明映射 \(f \mapsto f(\mathrm{D}_1, \ldots, \mathrm{D}_n)\) 是多项式代数 \(\mathbf{A}[\mathbf{X}_1, \ldots, \mathbf{X}_n] = \mathbf{E}\) 到 \(\mathbf{A}\) 上的自同态代数（即 \(\mathbf{A}\) -模 \(\mathbf{E}\) 的自同态代数）中的单同态（对 \(n\) 用归纳法）。
  \end{enumerate}
\end{enumerate}

\begin{enumerate}
\def\labelenumi{\alph{enumi})}
\setcounter{enumi}{1}
\tightlist
\item
  设 \(m\) 为整数 \(\geqslant 1\) ，使得 \(m: 1 = 0\) （在 \(\mathbf{A}\) 中）。证明 \(\mathrm{D}_i^m = 0\) 对每个偏导数 \(\mathrm{D}_i\) ( \(1 \leqslant i \leqslant n\) ) 在 \(\mathbf{A}[\mathbf{X}_1, \ldots, \mathbf{X}_n]\) 中成立。
\end{enumerate}

\begin{enumerate}
\def\labelenumi{\arabic{enumi})}
\setcounter{enumi}{4}
\tightlist
\item
  设 \(u \in \mathbf{A}[(\mathbf{X}_i)_{i \in I}]\) ， \(r\) 为整数 \(\geqslant 0\) ，并假设在 \(\mathbf{A}\) 中，对每个 \(n > 0\) ，关系 \(n\xi = 0\) 蕴含 \(\xi = 0\) 。如果我们有关系
\end{enumerate}

\[
\sum_ {i \in I} \left(\mathrm{D} _ {i} u\right) (\mathbf {X}) \mathbf {X} _ {i} = r u (\mathbf {X}),
\]

则 \(u\) 是 \(r\) 次齐次的。

\begin{enumerate}
\def\labelenumi{\arabic{enumi})}
\setcounter{enumi}{5}
\tightlist
\item
  设 K 为一个交换域， I 为一个集合， L 为环 K 上由集合 I 生成的自由结合代数（III，第448页，定义2），带有III第458页所定义的分次。我们按照IV第2页的方式定义 L 中元素的次数以及齐次元的概念。a) 证明 \(\deg(uv) = \deg(u) + \deg(v)\) 且 \(\deg(u + v) \leqslant \sup(\deg(u), \deg(v))\) （对 \(u, v \in L\) ）。
\end{enumerate}

\begin{enumerate}
\def\labelenumi{\alph{enumi})}
\setcounter{enumi}{1}
\item
  给定 \(u, v, q, q' \in L\) 满足 \(\deg(u) \geqslant \deg(v)\) 且 \(\deg(qu - q'v) < \deg(ku)\) ；证明存在 \(q'' \in L\) 使得 \(\deg(u - q''v) < \deg(u)\) 。
\item
  给定 \(u, v, q, q' \in L\) ，满足 \(\deg u \geqslant \deg v > \deg (qu - q'v)\) ；证明存在 \(q'' \in L\) 使得 \(\deg (u - q''v) < \deg v\) 。
\item
  给定 \(u, v, u', v' \in L\) ，齐次且满足 \(uv = u'v'\) 及 \(\deg v \geqslant \deg v' \geqslant 0\) ；存在 \(t \in L\) 使得 \(u' = ut, v = tv'\) 。
\item
  给定 \(u, v \in L\) ，齐次，满足 \(uv = vu\) ；存在 \(w \in L\) 、 \(a, b \in K\) 与 \(m, n \in K\) ，使得 \(u = aw^m\) , \(v = bw^n\) 。
\item
  设 \(u, v \in L\) 使得存在 \(a, b \in L\) （非零）满足 \(au = bv\) ；则存在 \(m, d \in L\) ，使得 \(u\) 与 \(v\) 公有的左倍式（相应地，右因式）是 \(m\) 的左倍式（相应地， \(d\) 的右因式）。
\end{enumerate}

